\documentclass[10pt,a4paper]{amsart}
\usepackage[margin=19mm]{geometry}
\usepackage{amsmath,amssymb,mathtools}
\usepackage[scr=rsfs,cal=euler]{mathalfa}
\usepackage{xfrac}
\usepackage[unicode,hidelinks,bookmarks=false]{hyperref}
\newcommand{\ie}{i.e.\ }
\newcommand{\cf}{cf.\ }
\newcommand{\resp}{respectively\ }
\newcommand{\Subsection}{\subsection}
\providecommand{\nobreakdash}{\nobreak\hbox{-}}
\setlength{\emergencystretch}{2em}
\multlinegap0pt
\usepackage{bm}
\usepackage{bbm}
\usepackage{dsfont}
\usepackage{xspace}
\usepackage{cancel}
\usepackage{mathtools}
\usepackage{amsthm}
\makeatletter
\@ifundefined{definition}{\newtheorem{definition}{Definition}}{}
\@ifundefined{lemma}{\newtheorem{lemma}[definition]{Lemma}}{}
\makeatother
\DeclareMathOperator{\GO}{GO}% the general orthogonal group
\DeclareMathOperator{\PGO}{PGO}% projective general orthogonal group
\title[Mixed identities for simple locally finite groups]{Mixed identities for simple locally finite groups}
\author{Henry Bradford, K{\i}van\c{c} Ersoy, Jakob Schneider,, Andreas Thom}
\date{Source: arXiv:2608.14537v1 (CC BY 4.0)}
\begin{document}
\maketitle

\noindent\emph{Source and licence.} Mixed identities for simple locally finite groups. Version arXiv:2608.14537v1 (CC BY 4.0), distributed under the Creative Commons Attribution 4.0 International licence, as recorded in the arXiv licence metadata for this exact version and shown on the abstract page https://arxiv.org/abs/2608.14537v1. Full rights notice: licenses/arxiv-2608.14537-CC-BY-4.0.txt.\par\smallskip

\noindent\textit{Editorial scope.} Counted unit: the Lemma whose source label is \texttt{None} in the article (source0.tex lines 1507--1509, source label \texttt{None}), delivered together with the article's own complete original proof of it (source0.tex lines 1511--1554, one \texttt{proof} environment of 44 lines). No mathematics is written, repaired or re-derived: statement and proof are the article's own text line for line. Required context retained uncounted: none. Every definition, display and label of those lines is retained unchanged. Omitted from this package and not redistributed: 1--1506, 1510--1510, 1555--2000. Nothing inside a retained excerpt is omitted or altered: every retained body is delivered exactly as the article's lines are written, so no omission receipt of any kind accompanies this package. Citations: the retained excerpts carry no citation command at all, so no bibliography entry of the article is transcribed and no reference is redistributed. Reference adaptation: none was needed and none was made; every reference inside the delivered unit resolves inside it, so no unresolved reference or citation is produced. Printed slips kept literal: no printed word, symbol, hypothesis or number of the counted statement or of its proof was altered. Layout and class: the article's own class is \texttt{amsproc} with packages including xargs, mathtools, todonotes, amsmath, amssymb, amsfonts, amsthm, hyperref, amsrefs, enumerate, geometry; the wrapper is the lane's pinned \texttt{amsart} editorial wrapper at 10pt on A4, which restates the theorem-like environments the retained text needs and adds the article's own macro definitions that the retained text needs (\texttt{\textbackslash GO}, \texttt{\textbackslash PGO}), with extra wrapper packages (bm, bbm, dsfont, xspace, cancel, mathtools). The delivered numbering is the unit's own. Assets: no retained span contains a figure environment, a graphics-inclusion command, a PDF-page inclusion, a table or a TikZ picture; no image file is copied into this package, the package declares no asset dependency and no image file is redistributed with it. Licence: version arXiv:2608.14537v1 of this article is distributed under the Creative Commons Attribution 4.0 International licence, as recorded in the arXiv licence metadata for this exact version and shown on the abstract page https://arxiv.org/abs/2608.14537v1. A full rights notice is delivered at licenses/arxiv-2608.14537-CC-BY-4.0.txt. Characters: every retained excerpt is pure ASCII, so the delivered engine, XeLaTeX with its default Latin Modern text font, reports no missing character.
\begin{lemma}
The group $\PGO_n^\bullet(q)$ has a mixed identity of length eight when $p\neq 2$.
\end{lemma}

\begin{proof}
Let $t\coloneqq 1_V+h$ be an \emph{Eichler transformation}, where 
$$
x.h=f(x,u)v-f(x,v)u
$$ 
for two linearly independent isotropic vectors $u\perp v$) and 
$$
x.l\coloneqq f(x,w)w
$$ 
a rank-one operator for $w\in V$ anisotropic such that the reflection 
$$
i\coloneqq r_w=1_V-2\frac{1}{f(w,w)}l
$$ 
lies in $\GO_n^\bullet(q)$. Let $g\in\GO_n^\bullet(q)$ be an arbitrary element. Define 
$$
c\coloneqq i^gti^gt=i^g(1_V+h)i^g(1_V+h)=1_V+i^ghi^g+h+i^ghi^gh.
$$
Regarding the last term in the previous equation, we see that (respecting $h^2=0$)
\begin{align}
i^ghi^gh&=\left(1-\frac{2}{f(w,w)}l\right)^gh\left(1-\frac{2}{f(w,w)}l\right)^gh\nonumber\\
&=\frac{4}{f(w,w)^2}l^ghl^gh-\frac{2}{f(w,w)}hl^gh\label{eq:help6}
\end{align}
Moreover, we compute:
\begin{align*}
x.l^ghl^g &=f(f(x,w.g)(w.g).h,w.g)w.g\\
&=f(x,w.g)f((w.g).h,w.g)w.g=0,
\end{align*}
since for any vector $v\in V$ we have 
\begin{align*}
f(v.(1+h),v)&=f(v,v.(1-h))=f(v,v)+f(v.h,v)\\&
=f(v,v)-f(v.h,v)
\end{align*}
implying that $f(v.h,v)=0$ (as $p\neq 2$). Take $v=w.g$ to complete the argument.
Whence we obtain in Equation~\eqref{eq:help6}
$$
i^ghi^gh=-\frac{2}{f(w,w)}hl^g h.
$$	
But then the critical term $i^ghi^gh$ in the expression of $c$ is invariant under conjugation with $i^g$:
\begin{align*}
i^g(i^ghi^gh)i^g&=\left(1-\frac{2}{f(w,w)}l^g\right)\left(-\frac{2}{f(w,w)}hl^g h\right)\left(1-\frac{2}{f(w,w)}l^g\right)\\
&=-\frac{2}{f(w,w)}hl^g h,
\end{align*}
since $l^ghl^g=0$. Then $w(g)=[c,i^g]=1$, and $w$ is of length eight and it descends to a mixed identity of $\PGO_n^\bullet(q)$. Thus we are done.
\end{proof}

\end{document}
